\documentclass[10pt,a4paper]{amsart}
\usepackage[margin=19mm]{geometry}
\usepackage{amsmath,amssymb,mathtools}
\usepackage[scr=rsfs,cal=euler]{mathalfa}
\usepackage{xfrac}
\usepackage[unicode,hidelinks,bookmarks=false]{hyperref}
\newcommand{\ie}{i.e.\ }
\newcommand{\cf}{cf.\ }
\newcommand{\resp}{respectively\ }
\newcommand{\Subsection}{\subsection}
\providecommand{\nobreakdash}{\nobreak\hbox{-}}
\setlength{\emergencystretch}{2em}
\multlinegap0pt
\usepackage{bm}
\usepackage{bbm}
\usepackage{dsfont}
\usepackage{xspace}
\usepackage{cancel}
\usepackage{mathtools}
\usepackage{amsthm}
\makeatletter
\@ifundefined{thm}{\newtheorem{thm}{Thm}}{}
\@ifundefined{lem}{\newtheorem{lem}[thm]{Lemma}}{}
\@ifundefined{prop}{\newtheorem{prop}[thm]{Proposition}}{}
\@ifundefined{rem}{\newtheorem{rem}{Remark}}{}
\makeatother
\makeatletter
\newcommand{\R}{\mathbb{R}}
\expandafter\ifx\csname case\endcsname\relax
\newenvironment{case}[1]{\par\noindent\textbf{Case 1.}\space#1}{}
\fi
\newcommand{\hr}{H^1_{0, rad}(\Omega)}
\expandafter\ifx\csname step\endcsname\relax
\newenvironment{step}[1]{\par\noindent\textbf{Step 1.}\space#1}{}
\fi
\makeatother
\title[Nonlocal problem with critical exponential nonlinearity of c]{Nonlocal problem with critical exponential nonlinearity of convolution type: A non-resonant case}
\author{Suman Kanungo, Pawan Kumar Mishra}
\date{Source: arXiv:2408.00654v2 (CC BY 4.0)}
\begin{document}
\maketitle

\noindent\emph{Source and licence.} Nonlocal problem with critical exponential nonlinearity of convolution type: A non-resonant case. Version arXiv:2408.00654v2 (CC BY 4.0), distributed under the Creative Commons Attribution 4.0 International licence, as recorded in the arXiv licence metadata for this exact version and shown on the abstract page https://arxiv.org/abs/2408.00654v2. Full rights notice: licenses/arxiv-2408.00654-CC-BY-4.0.txt.\par\smallskip

\noindent\textit{Editorial scope.} Counted unit: the Prop whose source label is \texttt{None} in the article (source0.tex lines 1243--1245, source label \texttt{None}), delivered together with the article's own complete original proof of it (source0.tex lines 1246--1389, one \texttt{proof} environment of 144 lines). No mathematics is written, repaired or re-derived: statement and proof are the article's own text line for line. Required context retained uncounted: J (lines 126--129); ceg (lines 144--150); AR (lines 237--245); H1 (lines 237--245); H4 (lines 237--245); d (lines 237--245); lem3.1 (lines 480--483); lem3.2 (lines 643--650); 3.14 (lines 653--655); rem3.1 (lines 754--757); 5.1 (lines 1159--1161); lem5.4 (lines 1174--1185); 5.3 (lines 1192--1196); 5.4 (lines 1232--1235). Every definition, display and label of those lines is retained unchanged. Omitted from this package and not redistributed: 1--125, 130--143, 151--236, 246--479, 484--642, 651--652, 656--753, 758--1158, 1162--1173, 1186--1191, 1197--1231, 1236--1242, 1390--1523. Nothing inside a retained excerpt is omitted or altered: every retained body is delivered exactly as the article's lines are written, so no omission receipt of any kind accompanies this package. Citations: the retained excerpts carry no citation command at all, so no bibliography entry of the article is transcribed and no reference is redistributed. Reference adaptation: none was needed and none was made; every reference inside the delivered unit resolves inside it, so no unresolved reference or citation is produced. Printed slips kept literal: no printed word, symbol, hypothesis or number of the counted statement or of its proof was altered. Layout and class: the article's own class is \texttt{amsart} with packages including amsmath, amssymb, latexsym, soul, cite, mathrsfs, amsfonts, float, xcolor, color, enumitem, graphicx, hyperref, babel; the wrapper is the lane's pinned \texttt{amsart} editorial wrapper at 10pt on A4, which restates the theorem-like environments the retained text needs and adds the article's own macro definitions that the retained text needs (\texttt{\textbackslash R}, \texttt{\textbackslash hr}), with extra wrapper packages (bm, bbm, dsfont, xspace, cancel, mathtools). The delivered numbering is the unit's own. Assets: no retained span contains a figure environment, a graphics-inclusion command, a PDF-page inclusion, a table or a TikZ picture; no image file is copied into this package, the package declares no asset dependency and no image file is redistributed with it. Licence: version arXiv:2408.00654v2 of this article is distributed under the Creative Commons Attribution 4.0 International licence, as recorded in the arXiv licence metadata for this exact version and shown on the abstract page https://arxiv.org/abs/2408.00654v2. A full rights notice is delivered at licenses/arxiv-2408.00654-CC-BY-4.0.txt. Characters: every retained excerpt is pure ASCII, so the delivered engine, XeLaTeX with its default Latin Modern text font, reports no missing character.
	\begin{equation}
		\label{J}
		J(u) = \frac{1}{2} \int\limits_\Omega |\nabla u|^2 \ dx - \frac{1}{2}\int\limits_\Omega \lambda u^2 \ dx -\frac{1}{2} \int\limits_\Omega \Bigg( \int\limits_\Omega \frac{Q(|y|)|F(u(y))}{|x-y|^{\mu}}dy \Bigg) Q(|x|)F(u(x)) \ dx.
	\end{equation}  

	\begin{equation} 
		\label{ceg}
		\lim_{|s| \rightarrow +\infty} \frac{|f(s)|}{e^{\alpha s^2}}=\begin{cases}
			0,& \forall \alpha > \alpha_0,\\
			+\infty, & \forall \alpha < \alpha_0.
		\end{cases}
	\end{equation}

	\begin{align} 
		\label{H1} \tag{H1} 
		&f\in C(\R), f(s)=0 \text{ for all } s\leq 0, f \text{ has critical exponential growth defined in \eqref{ceg}. }  \\
		&\text{  Moreover, } f(s)= o(s^{\frac{2 - \mu}{2}}). \nonumber\\
		&\text{ there exists }  K > 1 \text{ such that }  0< K F(s) \leq f(s) s , \text{ for all } s >0. \label{AR} \tag{H2}\\
		&\text{ there exist }  s_0 > 0 ,  M_0 > 0, \text{ and }  v \in (0, 1] \text{ such that }  0 < s^v F(s) \leq M_0 f(s),\text{ for all }\label{d} \tag{H3}\\ 
		&  s \geq s_0 .\nonumber\\
		&\liminf_{s\rightarrow + \infty} \frac{F(s)}{e^{\alpha_0s^2}} := \beta_0 >0. \tag{H4} \label{H4} 
	\end{align}

	\begin{lem}
		\label{lem3.1}
		Assume \eqref{H1}, \eqref{AR} holds and if $\{ u_n\} \subset \hr $ be a Palais-Smale sequence of $J$, then $\{u_n\}$ is a bounded sequence in $\hr$.
	\end{lem}

	\begin{lem} \label{lem3.2}
		Assume \eqref{H1} and \eqref{d} holds. If $\{u_n\}$ is a Palais-Smale sequence in $\hr$, then up to a subsequence, we have
		\begin{align*}
			\lim_{n \rightarrow \infty} &\int\limits_{\Omega} \Bigg( \int\limits_{\Omega} \frac{Q(|y|)F(u_n(y))}{|x-y|^\mu}dy \Bigg) Q(|x|)F(u_n(x))dx\\
			&=
			\int\limits_{\Omega} \Bigg( \int\limits_{\Omega} \frac{Q(|y|)F(u_0(y))}{|x-y|^\mu}dy \Bigg) Q(|x|)F(u_0(x))dx. 
		\end{align*}  
	\end{lem}

		\begin{align}
			\sup_n\int\limits_{\Omega} \Bigg( \int\limits_{\Omega} \frac{Q(|y|)F(u_n(y))}{|x-y|^\mu}dy \Bigg) Q(|x|)f(u_n(x))u_n(x)dx \leq C_0. \label{3.14}
		\end{align}

	\begin{rem}
		\label{rem3.1}
		Lemma \ref{lem3.2} holds even if we take $\{u_n\}$ to be a sequence converges weakly to $u_0$ in $\hr$ and satisfies \eqref{3.14}.
	\end{rem}

	\begin{equation} \label{5.1}
		W_n(x) = P_k(M_n(x))
	\end{equation} 

	\begin{lem}
		\label{lem5.4}
		Let $W_n$ be defined in \eqref{5.1}. Then the following estimates hold:
		\begin{enumerate}
			\item [(i)] $1 - \frac{A_0}{\log{n}} \leq \|W_n\|^2 \leq 1$.
			\item [(ii)] $W_n(x) \geq \left\{ 
			\begin{aligned}
				&\frac{-B_0}{\sqrt{\log{n}}},  &&\ \Omega\setminus B_{\frac{1}{n}}(0)\\
				& \frac{1}{\sqrt{2\pi}}\sqrt{\log{n}}- \frac{B_0}{\sqrt{\log{n}}}, && \ B_{\frac{1}{n}}(0)
			\end{aligned}\right. $
		\end{enumerate}
	\end{lem}

		\begin{align}
			\label{5.3}
			\|W_n\|^2 &\geq \|M_n\|^2 - A_0 \|(I-P_k)M_n\|_2^2 \nonumber\\
			&= 1 - A_0\|(I-P_k)M_n\|^2_2.
		\end{align}

	\begin{equation}
		\label{5.4}
		0<c(n) =\inf_{\gamma \in \Gamma_n} \max_{w\in \mathcal Q_n} J(\gamma(w)),
	\end{equation}

	\begin{prop}
		Let $c(n)$ be given as in \eqref{5.4} and assumptions \eqref{H1}-\eqref{H4} hold. Then for some $n$, $c(n)< \frac{(4-\mu)\pi}{2\alpha_0}\Big(1+\frac{2b_0}{4-\mu}\Big).$
	\end{prop}

	\begin{proof}Based on the definition of $c(n)$ and $id\in \Gamma_n$, it suffices to prove that $\max\{ J(v+sz_n): v \in H_{k,r}(\Omega), \|v\| \leq R, 0 \leq s \leq R\} < \frac{(4-\mu)\pi}{2\alpha_0}\Big(1+\frac{2b_0}{4-\mu}\Big)$. 
		On the contrary, let us assume this condition is not satisfied. Then for all $n \in \mathbb{N}$
		\[ \max\{ J(v+sz_n): v\in H_{k,r}, \|v\| \leq R, 0 \leq s \leq R\} \geq \frac{(4-\mu)\pi}{2\alpha_0}\Big(1+\frac{2b_0}{4-\mu}\Big).\]
		But note that for $s \geq R$, we have $J(v+sz_n)\leq0$. Hence, due to the compactness of $H_{k,r} \cap \overline{B_R}$, for each $n$, there exist $s_n>0$ and $v_n \in H_{k,r}$ such that
		\[ J(v_n+s_n z_n) = \max\{ J(v+sz_n): v\in H_{k,r}, \|v\| \leq R, 0 \leq s \leq R\}.\]
		This further implies that
		\begin{equation}
			\label{5.5}
			J(v_n+s_nz_n) \geq \frac{(4-\mu)\pi}{2\alpha_0}\Big(1+\frac{2b_0}{4-\mu}\Big).
		\end{equation}
		Let $u_n = v_n+s_nz_n$, then $J'(u_n)=0$ and we get
		\begin{equation}
			\label{5.6}
			\|u_n\|^2 -\lambda \|u_n\|^2_2 - \int\limits_\Omega \Bigg( \int\limits_\Omega \frac{Q(|y|)F(u_n(y))}{|x-y|^\mu}dy\Bigg) Q(|x|)f(u_n(x))u_n(x)dx =0.
		\end{equation}
		We complete the proof in the following two steps:\\
		Step 1. In this step, we prove that $\{v_n\}$ and $\{s_n\}$ are bounded sequences. 
		
		We have the following two distinct possibilities:
		\begin{enumerate}
			\item [(i)] $\frac{s_n}{\|v_n\|} \geq C_0,$ for some $C_0>0$ uniformly in $n$.
			\item [(ii)] $\frac{s_n}{\|v_n\|} \rightarrow 0$ in $\R$, up to a subsequence as $n \rightarrow \infty$. 
		\end{enumerate}
		Suppose $(i)$ holds. Then, the boundedness of $\{s_n\}$ implies the boundedness of $\{v_n\}$ as $\|v_n\| \leq \frac{s_n}{C_0}.$ Thus, it is sufficient to prove that the sequence $\{s_n\}$ is bounded. Condition $(i)$ implies that there exists a constant $C>0$ such that
		\begin{equation}
			\|u_n\| = \|v_n+s_nz_n\| \leq \|v_n\|+s_n\|z_n\| \leq \frac{s_n}{C_0} + s_n \leq \sqrt{C}s_n. \nonumber
		\end{equation}
		Using this in equation \eqref{5.6}, we get
		\begin{align}
			Cs_n^2 &\geq \int\limits\limits_\Omega \Big( \int\limits\limits_\Omega \frac{Q(|y|)F(u_n(y))}{|x-y|^\mu}dy\Big)Q(|x|)f(u_n(x)) u_n(x)dx \nonumber\\
			& \geq \int\limits\limits_{B_{\frac{1}{n}}} \Big( \int\limits\limits_{B_\frac{1}{n}} \frac{Q(|y|)F(u_n(y))}{|x-y|^\mu}dy\Big) Q(|x|)f(u_n(x))u_n(x)dx. \nonumber
		\end{align}
		The assumption \eqref{d} implies that, there exists $M_0>0$ and $s_0>0$ such that $f(s)s \geq M_0^{-1}s^{v+1} F(s)$ for all $s\geq s_0$. From \eqref{H4}, there exists $\varepsilon \in (0, \beta_0)$ such that $F(s) \geq (\beta_0-\varepsilon)e^{\alpha_0 s^2}$ for all $s \geq s_\varepsilon$. Choosing $\overline{s} = \max\{s_0, s_\varepsilon\}$, we have
		\begin{align}
			Cs_n^2 
			& \geq \frac{ (\beta_0-\varepsilon)^2}{M_0}\int\limits\limits_{B_{\frac{1}{n}} \cap \{ |u_n| \geq \overline{s}\}} \Big( \int\limits\limits_{B_\frac{1}{n} \cap \{ |u_n| \geq \overline{s}\}} \frac{|y|^{b_0}e^{\alpha_0 |u_n(y)|^2}}{|x-y|^\mu} dy\Big) |x|^{b_0}|u_n(x)|^{v+1} e^{\alpha_0 |u_n(x)|^2}dx. \label{5.7}
		\end{align}
		To estimate the above integral, we use Lemma \ref{lem5.4} $(i)$, $\|W_n(x)\|^2 \leq 1$ and from $(ii)$, in $B_{\frac{1}{n}}$ for $n$ large enough, we have
		\begin{align*}
			u_n(x) &= v_n(x) + s_n z_n(x) = v_n(x) + s_n \frac{W_n(x)}{\|W_n(x)\|}\\
			& \geq s_n \Big( \frac{1}{\sqrt{2\pi}}\sqrt{\log{n}} - \frac{B_0}{\log{n}}\Big) \Bigg( \frac{v_n(x)}{s_n\Big(\frac{1}{\sqrt{2\pi}}\sqrt{\log{n}}- \frac{B_0}{\log{n}}\Big)} + 1\Bigg)\\
			& \geq s_n \Big( \frac{1}{\sqrt{2\pi}}\sqrt{\log{n}} - \frac{B_0}{\log{n}}\Big) \Bigg( 1- \frac{ \|v_n\|_\infty}{s_n\Big(\frac{1}{\sqrt{2\pi}}\sqrt{\log{n}}- \frac{B_0}{\log{n}}\Big)} \Bigg).
		\end{align*}
		Since $H_{k,r}$ is a finite dimensional subspace, then there is a $C>0$ satisfying $\|v\|_\infty \leq C \|v\|$ for all $v\in H_{k,r}$. This implies that there exist $\delta', \delta'' \in (0,1)$, such that 
		\begin{align*}
			u_n(x)& \geq s_n \Big( \frac{1}{\sqrt{2\pi}}\sqrt{\log{n}} - \frac{B_0}{\log{n}}\Big)\Bigg( 1- \frac{ C\|v_n\|}{s_n\Big(\frac{1}{\sqrt{2\pi}}\sqrt{\log{n}}- \frac{B_0}{\log{n}}\Big)} \Bigg)\\
			& \geq s_n \Big( \frac{1}{\sqrt{2\pi}}\sqrt{\log{n}} - \frac{B_0}{\log{n}}\Big) \Bigg( 1- \frac{ C}{C_0\Big(\frac{1}{\sqrt{2\pi}}\sqrt{\log{n}}- \frac{B_0}{\log{n}}\Big)} \Bigg)\\
			& \geq \frac{s_n \sqrt{\log{n}}}{\sqrt{2\pi}} \Bigg( 1- \frac{ C}{C_0\Big(\frac{1}{\sqrt{2\pi}}\sqrt{\log{n}}- \frac{B_0}{\log{n}}\Big)} \Bigg)- \frac{B_0}{\log{n}} \Bigg( 1- \frac{ C}{C_0\Big(\frac{1}{\sqrt{2\pi}}\sqrt{\log{n}}- \frac{B_0}{\log{n}}\Big)} \Bigg)\\
			& \geq \frac{s_n \sqrt{\log{n}}}{\sqrt{2\pi}} (1-\delta') - \delta''.
		\end{align*}
		Choosing $ \delta \in (\delta', 1)$ such that $\frac{s_n \sqrt{\log{n}}}{\sqrt{2\pi}} (\delta-\delta') > \delta''$ for $n$ is large enough, then we have
		\begin{align*}
			u_n(x) &\geq \frac{s_n \sqrt{\log{n}}}{\sqrt{2\pi}} (1-\delta) + \frac{s_n \sqrt{\log{n}}}{\sqrt{2\pi}} (\delta-\delta') - \delta''\\
			& \geq \frac{s_n \sqrt{\log{n}}}{\sqrt{2\pi}} (1-\delta).
		\end{align*}
		
		Since $\frac{s_n \sqrt{\log{n}}}{\sqrt{2\pi}}( 1- \delta)$ is arbitrary large, we can take $\frac{s_n \sqrt{\log{n}}}{\sqrt{2\pi}}( 1- \delta) > \overline{s}$ for large $n$. Therefore, $B_{\frac{1}{n}} \subset \Big\{|u_n| \geq \frac{s_n \sqrt{\log{n}}}{\sqrt{2\pi}} ( 1- \delta) \Big\}$.
		From equation \eqref{5.7},
		\begin{align}
			Cs_n^2 &\geq \frac{ (\beta_0-\varepsilon)^2}{M_0 n^{2b_0}}\int\limits\limits_{B_{\frac{1}{n}}} \Big( \int\limits\limits_{B_\frac{1}{n}} \frac{e^{\frac{\alpha_0(1-\delta)^2 s_n^2 \log{n}}{2\pi}}}{|x-y|^\mu} dy\Big) \Big( \frac{s_n \sqrt{\log{n}}}{\sqrt{2\pi}} (1-\delta)\Big)^{v+1} e^{\frac{\alpha_0(1-\delta)^2 s_n^2 \log{n}}{2\pi}}dx \nonumber\\
			&\geq \frac{(\beta_0-\varepsilon)^2 e^{\frac{\alpha_0(1-\delta)^2 s_n^2 \log{n}}{\pi}}}{M_0 n^{2b_0}}\Big( \frac{s_n \sqrt{\log{n}}}{\sqrt{2\pi}} (1-\delta)\Big)^{v+1} \int\limits\limits_{B_{\frac{1}{n}}} \Big( \int\limits\limits_{B_\frac{1}{n}} \frac{1}{|x-y|^\mu} dy\Big) dx\nonumber  \\
			&\geq \frac{4\pi^2}{(2-\mu)(3-\mu)(4-\mu)} M_0^{-1} (\beta_0-\varepsilon)^2 e^{\frac{\alpha_0(1-\delta)^2 s_n^2 \log{n}}{\pi}}	\Big( \frac{s_n \sqrt{\log{n}}}{\sqrt{2\pi}} (1-\delta)\Big)^{v+1} n^{\mu-2b_0-4}. \nonumber
		\end{align}
		We get a constant $C_0= C_0(\mu, M_0, \overline{s}, \beta_0, \varepsilon)>0$ satisfying
		\begin{align}
			& s_n^2\geq C_0 (1-\delta)^{v+1} s_n^{v+1} e^{[\alpha_0 (1-\delta)^2 s_n^2 \pi^{-1}-(4+2b_0-\mu)]\log{n} + \frac{v+1}{2} \log{\log{n}}}. \label{5.8}
		\end{align}
		Since $ \frac{v+1}{2} \log{\log{n}} >0$, then
		\begin{equation}
			\label{5.9}
			s_n^2 \geq C_0 (1-\delta)^{v+1} s_n^{v+1} e^{[\alpha_0 (1-\delta)^2 s_n^2 \pi^{-1}-(4+2b_0-\mu)]\log{n}}. 
		\end{equation}
		By simple computations, we have 
		\begin{align}
			\label{5.10}
			\frac{(1-v)\log{s_n}}{s_n^2} \geq & \frac{\log{C_0(1-\delta)^{v+1}}}{s_n^2} + \left[\alpha_0 (1-\delta)^2 \pi^{-1}-\frac{(4+2b_0-\mu)}{s_n^2} \right]\log{n}.
		\end{align}
		If $s_n \rightarrow \infty$ in \eqref{5.10}, then we obtain a contradiction. Hence $\{s_n\}$ is a bounded sequence, so is $\{v_n\}$.
		
		Next, we assume that $(ii)$ occurs. Then, $s_n\leq \|v_n\|$, this gives that $\|u_n\|= \|v_n+ s_n z_n\| \leq \|v_n\|+ s_n \|z_n\|\leq 2 \|v_n\|$. This implis that if the sequence $\{v_n\}$ is bounded in $\hr$, then the sequence $\{s_n\}$ is bounded in $\R$. Therefore, our goal is to prove that $\{v_n\}$ is bounded in $\hr$. Let us assume that this is not true. Then, up to a subsequence $\|v_n\| \rightarrow \infty$. Let  $d= \sup\{ |x-y|: x,y \in \Omega\}$ denotes diameter of $\Omega$. From \eqref{5.6}, we have
		\begin{align}
			1 &\geq \frac{1}{\|u_n\|^2}\int\limits\limits_\Omega \Big( \int\limits\limits_\Omega \frac{Q(|y|)F(u_n(y))}{|x-y|^\mu}dy\Big) Q(|x|)f(u_n(x))u_n(x)\ dx \nonumber\\
			& \geq \frac{1}{4d^\mu \|v_n\|^2} \int\limits\limits_{ B_{\frac{1}{n}}} \Big( \int\limits\limits_{ B_{\frac{1}{n}}} |y|^{b_0}F(u_n(y))dy\Big) |x|^{b_0}f(u_n(x))u_n(x)dx. \nonumber 
		\end{align}
		Using \eqref{d} and \eqref{H4}, we can find $\varepsilon \in (0, \beta_0)$ and $\overline{s} = \max\{s_0, s_\varepsilon\}>0$ such that
		\begin{align}
			1 & \geq \frac{\overline{s}^{v+1}(\beta_0 - \varepsilon)^2}{4M_0d^\mu n^{2b_0} \|v_n\|^2} \int\limits\limits_{ B_{\frac{1}{n}} \cap \{|u_n| \geq \overline{s}\}} \Big( \int\limits\limits_{ B_{\frac{1}{n}} \cap \{|u_n| \geq \overline{s}\}} e^{\alpha_0 u_n^2(y)}dy\Big) e^{\alpha_0 u_n^2(x)}\ dx. \label{5.11}
		\end{align}
		Note that,
		\[ \frac{u_n}{\|v_n\|} \chi_{ B_{\frac{1}{n}} \cap \{|u_n| \geq \overline{s}\}} = \frac{v_n}{\|v_n\|} + \frac{s_n}{\|v_n\|}z_n - \frac{u_n}{\|v_n\|} \chi_{ B_{\frac{1}{n}} \cap \{|u_n| < \overline{s}\}}.\]
		Since $\frac{s_n}{\|v_n\|} \rightarrow 0$ in $\R$, $z_n \rightarrow 0$ a.e. in $\Omega$ and $\frac{u_n}{\|v_n\|} \rightarrow 0$ in $ B_{\frac{1}{n}} \cap \{|u_n| < \overline{s}\}$, then there exists $v_0 \in H_{k,r}(\Omega)$ such that 
		\begin{equation*}
			\frac{u_n(x)}{\|v_n\|} \chi_{ B_{\frac{1}{n}} \cap \{ |u_n| \geq \overline{s}\}} \rightarrow \overline{v} \text{ a.e. in } \Omega,
		\end{equation*}
		with $\frac{v_n}{\|v_n\|} \rightarrow \overline{v}$ and $\|\overline{v}\|=1$. From \eqref{5.11} and $\|v_n\| \to \infty$ as $n\to \infty$, 
		\begin{equation*}
			1\geq \frac{\overline{s}^{v+1}(\beta_0 - \varepsilon)^2}{4M_0d^\mu n^{2b_0}\|v_n\|^2} \int\limits\limits_\Omega \int\limits\limits_\Omega e^{\alpha_0 \|v_n\|^2 \Bigg( \frac{u_n(x)}{\|v_n\|} \chi_{\Omega \cap \{ |u_n| \geq \overline{s}\}}\Bigg)^2} e^{\alpha_0 \|v_n\|^2 \Bigg( \frac{u_n(x)}{\|v_n\|} \chi_{\Omega \cap \{ |u_n| \geq \overline{s}\}}\Bigg)^2}dx dy\rightarrow \infty, 
		\end{equation*}
		we get a contradiction. Hence, both $\{s_n\}$ and $\{v_n\}$ are bounded. \\
		Step 2. As we have proved that $\{s_n\}$ and $\{v_n\}$ are bounded, we can assume that up to a subsequence, there exists $v_0 \in \hr$ and $s_0\in \R$ such that $v_n \rightarrow v_0$ and $s_n \rightarrow s_0$. In this step, we prove that $v_0=0$ and $s_0^2 = \frac{(4-\mu)\pi}{\alpha_0}\Big( 1+\frac{2b_0}{4-\mu}\Big)$.
		
		Observe that since $M_n \rightharpoonup 0$, we have $ z_n \rightharpoonup 0$ and this implies $z_n \to 0$ a.e. in $\Omega$. Therefore, 
		\[u_n=v_n+s_n z_n \to v_0 \text{ a.e. in } \Omega.\]
		From \eqref{5.6} and Lemma \ref{lem3.1}, we have
		\begin{align*} \int\limits\limits_\Omega \Big( \int\limits\limits_\Omega \frac{Q(|y|)F(u_n(y))}{|x-y|^\mu}dy \Big) Q(|x|)f(u_n(x))u_n(x)dx = \|u_n\|^2 - \lambda\|u_n\|^2_2\leq  \|u_n\|^2 \leq C.
		\end{align*}
		By Remark \ref{rem3.1}, we have 
		\begin{equation} 
			\label{5.12}
			\int\limits\limits_\Omega \Big( \int\limits\limits_\Omega \frac{Q(|y|)F(u_n(y))}{|x-y|^\mu}dy \Big) Q(|x|)F(u_n(x)) dx \rightarrow \int\limits\limits_\Omega \Big( \int\limits\limits_\Omega \frac{Q(|y|)F(v_0(y))}{|x-y|^\mu}dy \Big) Q(|x|)F(v_0(x))dx.
		\end{equation}
		Since $v_n \rightarrow v_0$ in $\hr$, $\|z_n\|_2 \rightarrow 0$ and $s_n\to s_0$. Therefore, from \eqref{5.5}, \eqref{5.12} and the characterization of $\lambda_k$, we obtain
		\begin{align*}
			\frac{(4-\mu)\pi}{2\alpha_0} \Big( 1+\frac{2b_0}{4-\mu}\Big)&\leq \frac{1}{2} \|v_n\|^2 + \frac{s_n^2}{2} - \frac{\lambda}{2} \|v_n\|^2_2+ \frac{\lambda}{2} s_n^2 \|z_n\|_2^2\\
			&\quad- \frac{1}{2} \int\limits\limits_\Omega \Bigg( \int\limits\limits_\Omega \frac{Q(|y|)F(v_0(y))}{|x-y|^\mu}dy \Bigg) Q(|x|)F(v_0(x))dx\\
			& = \frac{1}{2} \|v_0\|^2 + \frac{s_0^2}{2} - \frac{\lambda}{2} \|v_0\|^2_2 \\
			&\quad- \frac{1}{2} \int\limits\limits_\Omega \Bigg( \int\limits\limits_\Omega \frac{Q(|y|)F(v_0(y))}{|x-y|^\mu}dy \Bigg) Q(|x|)F(v_0(x))dx.
		\end{align*}
		By the definition of the functional $J$ given in \eqref{J}, we have
		\begin{align}
			\label{6.1}
			J(v_0) + \frac{s_0^2}{2} \geq 	\frac{(4-\mu)\pi}{2\alpha_0} \Big( 1+\frac{2b_0}{4-\mu}\Big).
		\end{align}
		Since $J(v_0) \leq 0$ for all $v_0 \in H_{k,r}$, from Proposition \ref{5.3}. We have  $s_0^2 \geq \frac{(4-\mu)\pi}{\alpha_0} \Big( 1+\frac{2b_0}{4-\mu}\Big)$ and we know that $\|v_n\|\leq C$, for some $C>0$, this implies that $(ii)$ is not possible. Hence only $(i)$ is true, then from \eqref{5.9}
		\[ s_n^2 \geq C_0(1-\delta)^{v+1} s_n^{v+1} e^{[\alpha_0 (1-\delta)^2 s_n^2 \pi^{-1}-(4+2b_0-\mu)]\log{n}}.\]
		This implies that $s_0^2 \leq \frac{(4-\mu)\pi}{\alpha_0 (1-\delta)^2} \Big(1+\frac{2b_0}{4-\mu}\Big)$, otherwise we get a contradiction as $n\to \infty$. Taking $\delta$ very small close to 0, then we have $s_0^2 \leq \frac{(4-\mu)\pi}{\alpha_0}\Big(1+\frac{2b_0}{4-\mu}\Big)$. Hence $s_0^2 = \frac{(4-\mu)\pi}{\alpha_0}\Big(1+\frac{2b_0}{4-\mu}\Big)$.\\
		Next, we show that $v_0 \equiv 0$. From \eqref{6.1}, $J(v_0) \geq 0$ but $J(v_0) \leq 0$ for all $v_0 \in H_{k,r}$. Thus, $J(v_0)=0$. 
		
		By the characterization of $\lambda_k$, we have
		\begin{align*}
			0= J(v_0) & \leq \frac{1}{2} \Big(1 - \frac{\lambda}{\lambda_k}\Big)\|v_0\|^2 - \int\limits\limits_\Omega \Bigg( \int\limits\limits_\Omega \frac{Q(|y|)F(v_0(y))}{|x-y|^\mu}dy\Bigg) Q(|x|)F(v_0(x))dx.
		\end{align*}
		Using assumption \eqref{H1} and \eqref{AR}, we have $F(s)\geq 0$ for all $s\geq 0$. This together with $\lambda > \lambda_k$, we get
		\begin{align*}
			0 \leq \frac{1}{2} \Big(1 - \frac{\lambda}{\lambda_k}\Big)\|v_0\|^2 <0,
		\end{align*}
		Hence, $\|v_0\|=0$. This completes our claim.
		
		To complete the proof, note that from Step 2, up to a subsequence, we have $v_n \rightarrow 0$ in $H_{k,r}$ and $s_n \rightarrow s_0$. Then from equation \eqref{5.8}, letting $\delta \rightarrow 0^+$,
		\begin{align*}
			s_n^2 \geq C_0 s_n^{v+1}e^{[ \alpha_0 s_n^2 \pi^{-1} - (4+2b_0-\mu)]\log{n} + \frac{v+1}{2}\log{\log{n}}}.
		\end{align*}
		Now taking $n\rightarrow \infty$, then we get a contradiction. This completes the proof of the proposition.
	\end{proof}

\end{document}
